\chapter{The Rules of Differentiation, and Why They Work}

In the previous chapter you met the derivative as a \emph{limit}: the slope of a curve
at a point, the instantaneous rate of change, the number that the slopes of shrinking
secant lines close in on. That definition is the bedrock. But if you had to return to
the limit
\[
  f'(x) \;=\; \lim_{h\To 0}\frac{f(x+h)-f(x)}{h}
\]
every single time you wanted a derivative, calculus would be unusable. Differentiating
$x^7 - 4x^3 + \tfrac{1}{x}$ from scratch would take a page of algebra and a prayer that
no $h$ got left behind.

This chapter is about freedom from that drudgery. We will build a small toolkit of
\emph{rules} -- the power rule, the constant-multiple and sum rules, the product rule,
the quotient rule, and the crown jewel, the chain rule -- that let you differentiate
almost any function you can write down, quickly and mechanically. But this is a book
about \emph{understanding}, so we will never just hand you a rule. Every single one will
be \emph{derived} from the limit definition, and -- more importantly -- we will see
\emph{why} it has the shape it does. The product rule is not an arbitrary string of
symbols to memorize; it is the statement that a growing rectangle gains area on two
sides at once. The chain rule is not a mysterious ``multiply by the derivative of the
inside''; it is the plain fact that rates of change multiply when you link two processes
together.

By the end you will not only be able to differentiate fluently -- you will be able to
\emph{reconstruct} each rule if you ever forget it, because you will know where it comes
from.

\begin{focusbox}
The plan: each rule gets (1) a picture or piece of intuition, (2) a careful derivation
from the limit definition, (3) a \emph{why} box exposing the mechanism, and (4) worked
examples. The order is deliberate: the easy rules (power, constant multiple, sum) first,
then product and quotient, then the chain rule, which secretly powers everything that
comes after it in this book.
\end{focusbox}

\section{The Building Blocks: Two Derivatives We Already Know}
\label{ch11:building-blocks}

Before any rules, we need raw material. Two derivatives are so simple we can read them
straight off the limit definition, and everything else will be assembled from them.

\subsection{The derivative of a constant is zero}

Let $f(x) = c$, a constant -- the graph is a flat horizontal line at height $c$. A flat
line has slope $0$ everywhere, so we already expect $f'(x)=0$. The limit confirms it
instantly:
\[
  f'(x) = \lim_{h\To 0}\frac{f(x+h)-f(x)}{h}
        = \lim_{h\To 0}\frac{c - c}{h}
        = \lim_{h\To 0}\frac{0}{h}
        = \lim_{h\To 0} 0 = 0.
\]
The numerator is exactly zero for \emph{every} $h$, because a constant function does not
change. There is nothing accumulating, so there is no rate.

\begin{intuitionbox}
A constant function is a thing that never moves. Its rate of change is the rate at which
a parked car covers distance: zero. Keep this in your pocket -- it is why additive
constants (the ``$+C$'' you will meet again with integrals) vanish under
differentiation.
\end{intuitionbox}

\subsection{The derivative of $x$ is one}

Let $f(x)=x$. The graph is the $45^\circ$ line through the origin, slope $1$ everywhere.
Again the limit agrees:
\[
  f'(x) = \lim_{h\To 0}\frac{(x+h)-x}{h}
        = \lim_{h\To 0}\frac{h}{h}
        = \lim_{h\To 0} 1 = 1.
\]
When $x$ increases by $h$, the output $f(x)=x$ increases by exactly $h$. Output change
divided by input change is $h/h = 1$. The identity function passes change through
untouched, so its rate is $1$.

These two facts -- $\ddx{c}=0$ and $\ddx{x}=1$ -- are the atoms. Now let us learn how to
combine atoms.

\section{Constant Multiples and Sums: Differentiation is Linear}
\label{ch11:linearity}

\subsection{The constant-multiple rule}

Suppose you already know $f'$, and you scale the whole function by a constant $k$, forming
$g(x)=k\,f(x)$. What happens to the slope? Stretching a graph vertically by a factor of
$k$ stretches all its \emph{rises} by $k$ while leaving the \emph{runs} alone -- so every
slope is multiplied by $k$.

\begin{theorem}[Constant-multiple rule]\label{ch11:thm-constmult}
If $f$ is differentiable and $k$ is a constant, then $k f$ is differentiable and
\[
  \ddx{k\,f(x)} = k\,f'(x).
\]
\end{theorem}

\begin{proof}
Go to the definition and watch the constant ride along outside the difference quotient:
\[
  \big(kf\big)'(x)
  = \lim_{h\To 0}\frac{k f(x+h)-k f(x)}{h}
  = \lim_{h\To 0} k\cdot\frac{f(x+h)- f(x)}{h}
  = k\lim_{h\To 0}\frac{f(x+h)- f(x)}{h}
  = k\,f'(x).
\]
The single nontrivial move is pulling $k$ out of the limit, which is legal because a
constant factor always survives a limit unchanged.
\end{proof}

\subsection{The sum rule}

Now add two functions: $g(x)=f(x)+p(x)$. If $f$ is rising at rate $f'$ and $p$ is rising
at rate $p'$, then their sum rises at the combined rate $f'+p'$. Rates simply add, the
way two streams pouring into one tank add their flows.

\begin{theorem}[Sum (and difference) rule]\label{ch11:thm-sum}
If $f$ and $p$ are differentiable, then so is $f\pm p$, and
\[
  \ddx{f(x)\pm p(x)} = f'(x)\pm p'(x).
\]
\end{theorem}

\begin{proof}
Form the difference quotient and regroup the four terms so the $f$'s sit together and the
$p$'s sit together:
\[
  \big(f+p\big)'(x)
  = \lim_{h\To 0}\frac{[f(x+h)+p(x+h)] - [f(x)+p(x)]}{h}
  = \lim_{h\To 0}\!\left(\frac{f(x+h)-f(x)}{h} + \frac{p(x+h)-p(x)}{h}\right).
\]
The limit of a sum is the sum of the limits (each exists, by assumption), so this is
$f'(x)+p'(x)$. The subtraction case is identical with a minus sign.
\end{proof}

\begin{whybox}
Together, the constant-multiple and sum rules say the derivative is a \emph{linear
operator}: $\ddx{a f + b g} = a f' + b g'$. This is the single most useful structural
fact about differentiation. It is what lets you differentiate a long polynomial
term-by-term, never having to deal with the whole thing at once. The deep reason is that
the difference quotient itself is built from subtraction and division by $h$ -- both
linear operations -- and the limit preserves linearity. Differentiation inherits its
linearity from the linearity of ``measuring change.''
\end{whybox}

\begin{trapbox}
Linearity covers sums and constant multiples \emph{only}. It does \emph{not} extend to
products or quotients:
\[
  \ddx{f\cdot g}\;\ne\; f'\cdot g',
  \qquad
  \ddx{\frac{f}{g}}\;\ne\;\frac{f'}{g'}.
\]
A great many wrong answers in calculus come from wishing products behaved as nicely as
sums. They do not -- and the next two sections are devoted to what really happens.
\end{trapbox}

\section{The Power Rule}
\label{ch11:power-rule}

Here is the workhorse. For any constant exponent $n$,
\[
  \boxed{\;\ddx{x^n} = n\,x^{\,n-1}.\;}
\]
``Bring the exponent down in front, then knock the exponent down by one.'' You will use
this thousands of times. But where does it come from? We will prove it three ways, each
of which teaches something different: first for small integers by hand, then for all
positive integers by the binomial theorem, and finally we will explain why it keeps
working for negative and fractional exponents.

\subsection{Warm-up: $x^2$ and $x^3$ by hand}

\begin{example}[The derivative of $x^2$]\label{ch11:ex-x2}
Let $f(x)=x^2$. Then
\[
  f'(x)=\lim_{h\To 0}\frac{(x+h)^2 - x^2}{h}
       =\lim_{h\To 0}\frac{x^2 + 2xh + h^2 - x^2}{h}
       =\lim_{h\To 0}\frac{2xh + h^2}{h}
       =\lim_{h\To 0}(2x + h)
       = 2x.
\]
The $x^2$ cancels; the surviving terms all carry an $h$, so dividing by $h$ leaves
$2x+h$, and the leftover $h$ dies in the limit. Result: $\ddx{x^2}=2x = 2x^{2-1}$. The
pattern holds.
\end{example}

\begin{example}[The derivative of $x^3$]\label{ch11:ex-x3}
Now $f(x)=x^3$. Using $(x+h)^3 = x^3 + 3x^2 h + 3x h^2 + h^3$,
\[
  f'(x)=\lim_{h\To 0}\frac{3x^2 h + 3x h^2 + h^3}{h}
       =\lim_{h\To 0}\big(3x^2 + 3x h + h^2\big)
       = 3x^2 = 3x^{3-1}.
\]
Again the lone $x^3$ cancels, every surviving term has at least one $h$, and only the term with
a \emph{single} $h$ --- namely $3x^2 h$ --- contributes after we divide by $h$ and take
$h\To 0$. That single first-order term sets the instantaneous rate.
\end{example}

\begin{happeningbox}
Look at what survives. When you expand $(x+h)^n$ you get $x^n$ plus a term with exactly
one $h$, plus higher-order terms in $h^2, h^3,\dots$ The $x^n$ cancels the $-x^n$. After
dividing by $h$, the single-$h$ term loses its $h$ and becomes a clean constant in $h$; every
other term still carries at least one factor of $h$ and therefore vanishes as $h\To 0$.
\emph{The derivative reads off the coefficient of the linear term} ---
the portion of the change proportional to the nudge $h$. All higher powers
are second-order small, vanishing at the instant. This is the governing
mechanism of differential calculus: the derivative isolates the linear response of a function
to a microscopic step.
\end{happeningbox}

\subsection{All positive integers at once: the binomial theorem}

The hand calculations practically prove themselves once you name the pattern. The
binomial theorem expands
\[
  (x+h)^n
  = x^n + n x^{\,n-1} h + \binom{n}{2} x^{\,n-2} h^2 + \cdots + h^n .
\]
The first term is $x^n$, the second is $n x^{n-1}h$ (this is the one that matters), and
every later term carries $h^2$ or higher.

\begin{theorem}[Power rule, positive integers]\label{ch11:thm-power-int}
For every positive integer $n$, $\ddx{x^n} = n x^{\,n-1}$.
\end{theorem}

\begin{proof}
Subtract $x^n$, divide by $h$:
\[
  \frac{(x+h)^n - x^n}{h}
  = \frac{n x^{n-1}h + \binom{n}{2}x^{n-2}h^2 + \cdots + h^n}{h}
  = n x^{n-1} + \underbrace{\binom{n}{2}x^{n-2}h + \cdots + h^{n-1}}_{\text{every term has an } h}.
\]
As $h\To 0$ every term in the underbrace contains at least one factor of $h$ and so tends
to $0$. What remains is $n x^{n-1}$.
\end{proof}

\subsection{Why it survives to negative and fractional exponents}

The binomial proof only handles $n=1,2,3,\dots$ Yet the rule is true far more generally:
$\ddx{x^{-2}} = -2x^{-3}$ and $\ddx{x^{1/2}} = \tfrac12 x^{-1/2}$ both obey
``$n x^{n-1}$.'' We are not equipped to prove the full real-exponent case yet -- the
cleanest proof waits until we have the chain rule and the logarithm (next chapter), where
$x^n = e^{n\ln x}$ makes it a one-liner. But two special cases are within reach right now,
and they are worth doing because they showcase the algebra.

\begin{example}[A negative exponent: $1/x$]\label{ch11:ex-reciprocal}
Let $f(x)=x^{-1}=1/x$. Then
\[
  f'(x)=\lim_{h\To 0}\frac{1}{h}\!\left(\frac{1}{x+h}-\frac{1}{x}\right)
       =\lim_{h\To 0}\frac{1}{h}\cdot\frac{x-(x+h)}{x(x+h)}
       =\lim_{h\To 0}\frac{-h}{h\,x(x+h)}
       =\lim_{h\To 0}\frac{-1}{x(x+h)}
       = -\frac{1}{x^2}.
\]
That is $-x^{-2} = (-1)x^{-1-1}$: the power rule with $n=-1$, exactly as advertised.
\end{example}

\begin{example}[A fractional exponent: $\sqrt{x}$]\label{ch11:ex-sqrt}
Let $f(x)=\sqrt{x}=x^{1/2}$ (for $x>0$). The trick is to rationalize -- multiply by the
conjugate to clear the difference of square roots:
\[
  f'(x)=\lim_{h\To 0}\frac{\sqrt{x+h}-\sqrt{x}}{h}
       =\lim_{h\To 0}\frac{(\sqrt{x+h}-\sqrt{x})(\sqrt{x+h}+\sqrt{x})}{h(\sqrt{x+h}+\sqrt{x})}.
\]
The numerator collapses to $(x+h)-x = h$, which cancels the $h$ below:
\[
  f'(x)=\lim_{h\To 0}\frac{1}{\sqrt{x+h}+\sqrt{x}}
       = \frac{1}{2\sqrt{x}}
       = \tfrac12\,x^{-1/2}.
\]
And $\tfrac12 x^{-1/2} = \tfrac12 x^{\,1/2 - 1}$ -- the power rule again, now with
$n=\tfrac12$.
\end{example}

\begin{whybox}
Why does ``bring down the exponent'' work? Geometrically, think of $x^n$ as the
\emph{volume of an $n$-dimensional cube} of side $x$. Nudge the side from $x$ to $x+h$.
The cube gains a thin slab on each of its $n$ faces, each slab of area $x^{n-1}$ and
thickness $h$, for a leading gain of $n\cdot x^{n-1}\cdot h$. (The edges and corners where
slabs overlap involve $h^2$ and smaller -- negligible at an instant.) Divide that gain by
$h$ and you get $n x^{n-1}$. The $n$ in front is literally counting the $n$ faces of the
cube that grow. We will see this same ``thicken each face'' picture again, in two
dimensions, when we prove the product rule.
\end{whybox}

\begin{figure}[H]
\centering
\begin{tikzpicture}
\begin{axis}[
  width=11cm, height=7.2cm,
  xlabel={$x$}, ylabel={$y$},
  domain=-1.7:1.7, samples=120,
  axis lines=middle,
  xmin=-1.9, xmax=1.9, ymin=-2.4, ymax=3.2,
  legend style={at={(0.5,1.05)}, anchor=south, legend columns=2, font=\footnotesize},
  legend cell align=left,
  enlargelimits=false,
]
\addplot[curveblue, very thick] {x^2};
\addlegendentry{$y=x^2$}
\addplot[tangentred, thick, samples=2, domain=-0.2:2.2] {2*1*(x-1)+1};
\addlegendentry{tangent at $x=1$, slope $2x=2$}
\addplot[structurepurple, thick, dashed, samples=2, domain=-2.2:0.2] {2*(-1)*(x+1)+1};
\addlegendentry{tangent at $x=-1$, slope $-2$}
\addplot[only marks, mark=*, black] coordinates {(1,1) (-1,1)};
\end{axis}
\end{tikzpicture}
\caption{The power rule in action for $y=x^2$. At $x=1$ the predicted slope is $2x=2$;
at $x=-1$ it is $-2$. The tangent lines have exactly those slopes. Differentiation reads
the steepness directly off the formula $2x$.}
\label{ch11:fig-power}
\end{figure}

\subsection{Polynomials: linearity plus the power rule}

The power rule plus linearity (Section~\ref{ch11:linearity}) differentiate \emph{any}
polynomial on sight, term by term.

\begin{example}[A polynomial]\label{ch11:ex-poly}
Differentiate $p(x) = 4x^5 - 7x^3 + 2x^2 - 9x + 6$.

Work term by term. The exponent comes down, drops by one, and constants multiply through;
the lone constant $6$ dies:
\[
  p'(x) = 4\cdot 5x^4 - 7\cdot 3x^2 + 2\cdot 2x - 9\cdot 1 + 0
        = 20x^4 - 21x^2 + 4x - 9.
\]
\end{example}

\begin{example}[Rewrite, then differentiate]\label{ch11:ex-rewrite}
Differentiate $g(x) = \dfrac{3}{x^2} + 5\sqrt{x} - \dfrac{x^4}{2}$.

First rewrite every piece as a power of $x$ so the power rule can grip it:
\[
  g(x) = 3x^{-2} + 5x^{1/2} - \tfrac12 x^4 .
\]
Now differentiate each power:
\[
  g'(x) = 3(-2)x^{-3} + 5\cdot\tfrac12 x^{-1/2} - \tfrac12\cdot 4x^3
        = -6x^{-3} + \tfrac{5}{2}x^{-1/2} - 2x^3
        = -\frac{6}{x^3} + \frac{5}{2\sqrt{x}} - 2x^3 .
\]
The lesson: rewriting roots and fractions as exponents is half the battle.
\end{example}

\begin{labbox}
Differentiate $h(x) = x^7 - 3x^4 + \dfrac{1}{x^3} - \sqrt[3]{x}$ by hand. (Hint: rewrite
the last two terms as $x^{-3}$ and $x^{1/3}$.) You should get
$h'(x) = 7x^6 - 12x^3 - 3x^{-4} - \tfrac13 x^{-2/3}$. Then check one value, say
$x=1$, by comparing to a numerical slope $\dfrac{h(1.001)-h(0.999)}{0.002}$.
\end{labbox}

\section{The Product Rule}
\label{ch11:product-rule}

We come to the first rule that genuinely surprises people. The naive guess --
``the derivative of a product is the product of the derivatives'' -- is \emph{wrong}, and
a single example demolishes it.

\begin{trapbox}
Take $f(x)=x$ and $g(x)=x$, so $f(x)g(x)=x^2$. The correct derivative is $2x$. But
$f'(x)g'(x) = 1\cdot 1 = 1 \ne 2x$. The product of the derivatives is not the derivative
of the product. Something else is going on.
\end{trapbox}

\subsection{The picture: a growing rectangle}

One picture makes the product rule obvious. Imagine a rectangle whose width is
$f(x)$ and whose height is $g(x)$. Its area is the product $A(x)=f(x)g(x)$. Now nudge $x$
forward by a tiny $h$. The width grows by $\Delta f = f(x+h)-f(x)$ and the height grows by
$\Delta g = g(x+h)-g(x)$. How much new area appears?

\begin{figure}[H]
\centering
\begin{tikzpicture}[scale=0.95]
  % Panel 1: Geometric Expansion
  \begin{scope}[shift={(0,0)}]
    \node[font=\footnotesize\scshape, text=gray!70!black] at (2.5, 4.2) {Geometric Area Expansion};

    % original rectangle
    \fill[sourceblue!15, rounded corners=2pt] (0,0) rectangle (3.6,2.4);
    \draw[thick, sourceblue!70] (0,0) rectangle (3.6,2.4);
    \node[sourceblue!60!black, font=\bfseries] at (1.8,1.2) {$f(x)\cdot g(x)$};

    % right strip (width grows): g * Delta f
    \fill[targetgreen!20, rounded corners=2pt] (3.6,0) rectangle (5.0,2.4);
    \draw[thick, targetgreen!70] (3.6,0) rectangle (5.0,2.4);
    \node[align=center, targetgreen!60!black, font=\small] at (4.3,1.2) {$g\,\Delta f$};

    % top strip (height grows): f * Delta g
    \fill[quotientorange!20, rounded corners=2pt] (0,2.4) rectangle (3.6,3.4);
    \draw[thick, quotientorange!70] (0,2.4) rectangle (3.6,3.4);
    \node[quotientorange!60!black, font=\small] at (1.8,2.9) {$f\,\Delta g$};

    % corner: Delta f * Delta g
    \fill[trapred!25, rounded corners=2pt] (3.6,2.4) rectangle (5.0,3.4);
    \draw[thick, trapred!70] (3.6,2.4) rectangle (5.0,3.4);
    \node[trapred!70!black, font=\scriptsize] at (4.3,2.9) {$\Delta f\Delta g$};

    % Dimension markers
    \draw[<->, >=stealth, gray!70!black] (0,-0.4) -- (3.6,-0.4) node[midway, below, font=\small]{$f(x)$};
    \draw[<->, >=stealth, targetgreen!70!black] (3.6,-0.4) -- (5.0,-0.4) node[midway, below, font=\small]{$\Delta f$};
    \draw[<->, >=stealth, gray!70!black] (-0.4,0) -- (-0.4,2.4) node[midway, left, font=\small]{$g(x)$};
    \draw[<->, >=stealth, quotientorange!70!black] (-0.4,2.4) -- (-0.4,3.4) node[midway, left, font=\small]{$\Delta g$};
  \end{scope}

  % Divider corridor
  \draw[gray!30, dashed] (5.8, -0.6) -- (5.8, 4.4);

  % Panel 2: Rate Breakdown Ladder (art_costs lane architecture)
  \begin{scope}[shift={(6.6,0)}]
    \node[font=\footnotesize\scshape, text=gray!70!black] at (4.1, 4.2) {Rate Decomposition};

    % Lane 1: Width growth
    \draw[targetgreen!60, fill=targetgreen!12, rounded corners=2pt] (0,2.6) rectangle (8.2,3.5);
    \node[anchor=west, font=\footnotesize] at (0.2, 3.05) {\textbf{Width growth}:};
    \node[anchor=east, font=\footnotesize] at (8.0, 3.05) {$g(x)\,\dfrac{\Delta f}{h} \longrightarrow g(x)\,f'(x)$};

    % Lane 2: Height growth
    \draw[quotientorange!60, fill=quotientorange!12, rounded corners=2pt] (0,1.4) rectangle (8.2,2.3);
    \node[anchor=west, font=\footnotesize] at (0.2, 1.85) {\textbf{Height growth}:};
    \node[anchor=east, font=\footnotesize] at (8.0, 1.85) {$f(x)\,\dfrac{\Delta g}{h} \longrightarrow f(x)\,g'(x)$};

    % Lane 3: Second-order corner
    \draw[trapred!60, fill=trapred!10, rounded corners=2pt] (0,0.2) rectangle (8.2,1.1);
    \node[anchor=west, font=\footnotesize] at (0.2, 0.65) {\textbf{Corner interaction}:};
    \node[anchor=east, font=\footnotesize] at (8.0, 0.65) {$\Delta f\cdot\dfrac{\Delta g}{h} \longrightarrow 0\cdot g'(x) = 0$};

    % Total summary bar
    \draw[sourceblue!80, fill=sourceblue!15, rounded corners=2pt, very thick] (0,-0.8) rectangle (8.2,-0.1);
    \node[anchor=west, font=\footnotesize\bfseries] at (0.2, -0.45) {Total rate:};
    \node[anchor=east, font=\footnotesize\bfseries] at (8.0, -0.45) {$\ddx{f(x)\,g(x)} = f'(x)g(x) + f(x)g'(x)$};
  \end{scope}
\end{tikzpicture}
\caption{Why the product rule has two terms. Left: Growing both dimensions of a rectangle
creates a tall strip $g\,\Delta f$, a wide strip $f\,\Delta g$, and a microscopic corner
$\Delta f\,\Delta g$. Right: Dividing each piece by $h$ and taking the limit $h\to 0$ reveals that
the two edge strips produce finite rates $g f'$ and $f g'$, while the second-order corner vanishes
because $\Delta f \to 0$.}
\label{ch11:fig-product}
\end{figure}

The new area comes in three pieces (Figure~\ref{ch11:fig-product}):
\begin{itemize}
  \item a \emph{tall strip} along the right edge, of height $g$ and width $\Delta f$:
        area $g\,\Delta f$;
  \item a \emph{wide strip} along the top edge, of width $f$ and height $\Delta g$:
        area $f\,\Delta g$;
  \item a \emph{tiny corner} rectangle of dimensions $\Delta f$ by $\Delta g$:
        area $\Delta f\,\Delta g$.
\end{itemize}
So the change in area is
\[
  \Delta A = g\,\Delta f + f\,\Delta g + \Delta f\,\Delta g .
\]
Divide by $h$ to get the rate, and watch the corner term:
\[
  \frac{\Delta A}{h} = g\,\frac{\Delta f}{h} + f\,\frac{\Delta g}{h}
                       + \Delta f\cdot\frac{\Delta g}{h}.
\]
As $h\To 0$, $\Delta f/h \To f'$ and $\Delta g/h \To g'$. The first two terms become
$g f'$ and $f g'$. The corner term is $\Delta f\cdot(\Delta g/h)$; the factor $\Delta g/h$
tends to $g'$ (finite), but $\Delta f \To 0$, so the whole corner vanishes.
The geometric insight is straightforward: \emph{the corner is the product of two small quantities,
and therefore shrinks at second order}. It contributes nothing to the instantaneous rate of change.

\begin{theorem}[Product rule]\label{ch11:thm-product}
If $f$ and $g$ are differentiable, then so is $fg$, and
\[
  \ddx{f(x)\,g(x)} = f'(x)\,g(x) + f(x)\,g'(x).
\]
\end{theorem}

\begin{proof}
The picture is correct, but let us pin it down with the clean algebraic trick that makes
the corner term appear and disappear: \emph{add and subtract the same mixed term}. Start
from the difference quotient and insert $-f(x+h)g(x)+f(x+h)g(x)$ (a clever zero):
\[
  \frac{f(x+h)g(x+h) - f(x)g(x)}{h}
  = \frac{f(x+h)g(x+h) - f(x+h)g(x) + f(x+h)g(x) - f(x)g(x)}{h}.
\]
Group the first two and last two terms, factoring out the common piece in each:
\[
  = f(x+h)\,\frac{g(x+h)-g(x)}{h} \;+\; g(x)\,\frac{f(x+h)-f(x)}{h}.
\]
Now take $h\To 0$ term by term. In the first term, $f(x+h)\To f(x)$ (because $f$, being
differentiable, is continuous) and $\dfrac{g(x+h)-g(x)}{h}\To g'(x)$. In the second term,
$g(x)$ is constant in $h$ and $\dfrac{f(x+h)-f(x)}{h}\To f'(x)$. Hence the limit is
\[
  f(x)\,g'(x) + g(x)\,f'(x),
\]
which is the product rule.
\end{proof}

\begin{whybox}
Why \emph{two} terms? Because a product depends on \emph{both} of its factors, and at any
instant both factors are changing at once. The rectangle picture makes this vivid: area
is gained on two sides simultaneously. The term $f'g$ is the contribution from the width
growing while the height holds; $fg'$ is the contribution from the height growing while
the width holds. You can think of it as ``hold one factor still, differentiate the other,
then swap and add.'' The reason you \emph{add} them (rather than, say, multiply) is that
to first order the two growths are independent -- their interaction, the corner
$\Delta f\,\Delta g$, is second-order small and drops out. This ``vary one factor at a
time and sum'' pattern is exactly how the product rule generalizes: for three factors,
\[
  (fgh)' = f'gh + fg'h + fgh',
\]
one term per factor, each differentiating a single factor while freezing the others.
\end{whybox}

\begin{example}[A first product]\label{ch11:ex-prod1}
Differentiate $y = (x^2+1)(x^3 - 4x)$.

Let $f = x^2+1$ and $g = x^3 - 4x$, so $f' = 2x$ and $g' = 3x^2 - 4$. Then
\[
  y' = f'g + fg' = (2x)(x^3-4x) + (x^2+1)(3x^2-4).
\]
Expand to simplify:
\[
  y' = (2x^4 - 8x^2) + (3x^4 - 4x^2 + 3x^2 - 4) = 5x^4 - 9x^2 - 4 .
\]
\emph{Sanity check.} We could instead multiply out first:
$y = x^5 - 4x^3 + x^3 - 4x = x^5 - 3x^3 - 4x$, whose derivative is
$5x^4 - 9x^2 - 4$. Same answer. The product rule and ``expand first'' must agree -- and
here expanding was about as fast. The product rule earns its keep when expanding is
ugly or impossible, as with $x^2\sin x$ or $e^x\ln x$, which you cannot multiply out.
\end{example}

\begin{example}[When you cannot expand]\label{ch11:ex-prod2}
Differentiate $y = x^2 \sqrt{x+1}$. Write $g=(x+1)^{1/2}$, so (peeking ahead to the chain
rule, or just by the power rule on $x+1$) $g' = \tfrac12 (x+1)^{-1/2}$. With $f=x^2$,
$f'=2x$:
\[
  y' = 2x\,\sqrt{x+1} + x^2\cdot\frac{1}{2\sqrt{x+1}}
     = 2x\sqrt{x+1} + \frac{x^2}{2\sqrt{x+1}} .
\]
There is no way to ``expand'' a square root away; the product rule is essential here.
\end{example}

\section{The Quotient Rule}
\label{ch11:quotient-rule}

Division is just multiplication in disguise, so we can squeeze the quotient rule out of
the product rule. Suppose $q(x) = \dfrac{f(x)}{g(x)}$ (with $g(x)\ne 0$). The cleanest
derivation clears the denominator first.

\subsection{Deriving it from the product rule}

Write $f = q\cdot g$. Differentiate both sides with the product rule:
\[
  f' = q' g + q g'.
\]
Now solve for the unknown $q'$:
\[
  q' = \frac{f' - q g'}{g} = \frac{f' - \dfrac{f}{g}\,g'}{g}
     = \frac{f' g - f g'}{g^2},
\]
where in the last step we substituted $q=f/g$ and multiplied top and bottom by $g$.

\begin{theorem}[Quotient rule]\label{ch11:thm-quotient}
If $f$ and $g$ are differentiable and $g(x)\ne 0$, then $f/g$ is differentiable and
\[
  \ddx{\frac{f(x)}{g(x)}}
  = \frac{f'(x)\,g(x) - f(x)\,g'(x)}{\big(g(x)\big)^2}.
\]
\end{theorem}

\begin{remark}
The derivation above quietly assumed $q$ was differentiable so that we could apply the
product rule to $f=qg$. A fully airtight proof goes back to the limit definition (combine
the fractions over the common denominator $g(x+h)g(x)$, add and subtract a mixed term just
as we did for the product rule). The shortcut via $f=qg$ is the way to \emph{remember}
the rule; the limit proof is the way to \emph{certify} it. Both give the same formula.
\end{remark}

\begin{whybox}
The quotient rule's strange-looking minus sign and squared denominator are not arbitrary.
Think of $f/g$ as $f$ times $1/g$. Increasing $f$ pushes the quotient \emph{up}: that is
the $+f' g$ part. Increasing $g$ pushes the quotient \emph{down} -- a bigger denominator
means a smaller fraction -- which is exactly why $g'$ enters with a \emph{minus}. The
$g^2$ in the denominator is the price of the reciprocal: as we computed back in
Example~\ref{ch11:ex-reciprocal}, $\ddx{1/g}$ carries a $1/g^2$, and the quotient rule
inherits it. Order matters in the numerator: it is $f'g - fg'$, ``derivative of top times
bottom, minus top times derivative of bottom.'' Swapping the two terms flips the sign and
gives the wrong answer.
\end{whybox}

\begin{figure}[H]
\centering
\begin{tikzpicture}[scale=0.92]
  % Opposing rates on a fraction q = f/g
  % Left Panel: Numerator variation
  \begin{scope}[shift={(0,0)}]
    \node[font=\footnotesize\scshape, text=targetgreen!70!black] at (3.3, 3.4) {Numerator Growth ($f \to f + \Delta f$)};
    \draw[targetgreen!60, fill=targetgreen!10, rounded corners=2pt] (0,0) rectangle (6.6,2.8);

    \node[anchor=north west, font=\small] at (0.2, 2.6) {Holding denominator $g$ fixed:};
    \node[font=\small] at (3.3, 1.7) {$\dfrac{f + \Delta f}{g} - \dfrac{f}{g} \;=\; \dfrac{\Delta f}{g}$};
    \draw[->, >=stealth, thick, targetgreen!80!black] (3.3, 1.2) -- (3.3, 0.6)
      node[midway, right, font=\footnotesize] {divide by $h$, $h\to 0$};
    \node[font=\small\bfseries, targetgreen!70!black] at (3.3, 0.3) {$+\,\dfrac{f'(x)}{g(x)} \;=\; +\,\dfrac{f'(x)\,g(x)}{\big(g(x)\big)^2}$};
  \end{scope}

  % Right Panel: Denominator variation
  \begin{scope}[shift={(7.4,0)}]
    \node[font=\footnotesize\scshape, text=quotientorange!80!black] at (3.3, 3.4) {Denominator Growth ($g \to g + \Delta g$)};
    \draw[quotientorange!60, fill=quotientorange!10, rounded corners=2pt] (0,0) rectangle (6.6,2.8);

    \node[anchor=north west, font=\small] at (0.2, 2.6) {Holding numerator $f$ fixed:};
    \node[font=\small] at (3.3, 1.7) {$\dfrac{f}{g + \Delta g} - \dfrac{f}{g} \;=\; f\cdot\dfrac{-\Delta g}{g(g+\Delta g)}$};
    \draw[->, >=stealth, thick, quotientorange!80!black] (3.3, 1.2) -- (3.3, 0.6)
      node[midway, right, font=\footnotesize] {divide by $h$, $h\to 0$};
    \node[font=\small\bfseries, quotientorange!80!black] at (3.3, 0.3) {$-\,\dfrac{f(x)\,g'(x)}{\big(g(x)\big)^2}$};
  \end{scope}

  % Bottom Panel: Combined result
  \draw[sourceblue!80, fill=sourceblue!12, rounded corners=2pt, very thick] (0,-1.3) rectangle (14.0,-0.3);
  \node[font=\small] at (7.0, -0.8) {\textbf{Combined Net Rate}: $\displaystyle\ddx{\frac{f(x)}{g(x)}} \;=\; \frac{f'(x)\,g(x) - f(x)\,g'(x)}{\big(g(x)\big)^2}$};
\end{tikzpicture}
\caption{The mechanics of the quotient rule (Figure~\ref{ch11:fig-quotient}). The ratio $f/g$ responds to
two competing effects: numerator growth lifts the quotient at rate $+f'/g = +(f'g)/g^2$, while denominator
growth pulls the quotient down at rate $-(fg')/g^2$. Superposing both linear responses yields the
difference $f'g - fg'$ over the squared denominator $g^2$.}
\label{ch11:fig-quotient}
\end{figure}

\begin{trapbox}
The numerator of the quotient rule is \emph{not} symmetric: $f'g - fg' \ne fg' - f'g$. A
mnemonic that prevents the classic sign error: ``\,low dee-high minus high dee-low, over
the square of what's below.\,'' Here ``high'' is the numerator $f$, ``low'' is the
denominator $g$, and ``dee'' means ``derivative of.'' So:
$\dfrac{\text{low}\cdot f' - \text{high}\cdot g'}{\text{low}^2}
= \dfrac{g f' - f g'}{g^2}$. Same thing, written to keep the order straight.
\end{trapbox}

\begin{example}[A rational function]\label{ch11:ex-quot1}
Differentiate $q(x) = \dfrac{x^2 - 1}{x^2 + 1}$.

Top $f = x^2-1$, $f' = 2x$. Bottom $g = x^2+1$, $g' = 2x$. Then
\[
  q'(x) = \frac{f'g - fg'}{g^2}
        = \frac{2x(x^2+1) - (x^2-1)(2x)}{(x^2+1)^2}.
\]
Expand the numerator:
\[
  2x(x^2+1) - 2x(x^2-1) = 2x\big[(x^2+1)-(x^2-1)\big] = 2x\cdot 2 = 4x.
\]
So $q'(x) = \dfrac{4x}{(x^2+1)^2}$. Notice the numerator is clean while the denominator
stays squared -- a typical quotient-rule outcome.
\end{example}

\begin{example}[Deriving the reciprocal rule]\label{ch11:ex-quot2}
Take $f = 1$ (so $f' = 0$) in the quotient rule:
\[
  \ddx{\frac{1}{g}} = \frac{0\cdot g - 1\cdot g'}{g^2} = -\frac{g'}{g^2}.
\]
This \emph{reciprocal rule} is worth memorizing on its own. As a check, with $g=x$ it
gives $\ddx{1/x} = -1/x^2$, matching Example~\ref{ch11:ex-reciprocal}.
\end{example}

\begin{historybox}
The rules of this chapter were forged in the dramatic collision between two radically different mathematical worldviews during the 1670s and 1680s. Isaac Newton, working in Cambridge during the plague years of 1665--1666, conceived of calculus kinematically. For Newton, curves were physical trajectories traced out by moving points, fluent quantities changing continuously through time, and rates of change were \emph{fluxions} ($\dot{x}, \dot{y}$). Newton possessed the product rule and chain rule early, but he viewed them through the mechanics of motion and guarded his manuscripts privately, sharing them only with an inner circle of English correspondents.

Gottfried Wilhelm Leibniz, working independently in Paris and Hanover, approached the problem not from kinematics, but from combinatorics and symbolic logic. In October 1684, Leibniz published the very first paper on the differential calculus in the Leipzig journal \emph{Acta Eruditorum}: \emph{Nova Methodus pro Maximis et Minimis, itemque Tangentibus...} (``A New Method for Maxima and Minima, as well as Tangents, which is neither impeded by fractional nor irrational quantities...''). At just six pages, Leibniz's paper changed the mathematical world because of its notation and operational perspective. Leibniz introduced the differential symbols $\dd x$ and $\dd y$, and treated the product and quotient rules as principles of \emph{geometric algebra}. 

To Leibniz, the product $u \cdot v$ was the area of a rectangle with sides $u$ and $v$. If each side is given an infinitesimal increment $\dd u$ and $\dd v$, the new area is $(u + \dd u)(v + \dd v) = uv + u\,\dd v + v\,\dd u + \dd u\,\dd v$. By discarding the second-order corner $\dd u\,\dd v$ as incomparably smaller than the first-order strips, Leibniz deduced $\dd(uv) = u\,\dd v + v\,\dd u$ directly. To find the differential of a quotient $q = u/v$, he inverted the geometry: writing $u = qv$, he applied the product rule to get $\dd u = q\,\dd v + v\,\dd q$, and solved algebraically for $\dd q = (\dd u - q\,\dd v)/v = (v\,\dd u - u\,\dd v)/v^2$.

This operational elegance ignited a bitter priority dispute between the Royal Society and continental mathematicians that lasted decades. It also invited fierce philosophical attack: Bishop George Berkeley, in his 1734 critique \emph{The Analyst}, famously mocked these discarded differentials as ``the ghosts of departed quantities''---vanishing increments that were treated as non-zero when dividing, and zero when discarding corners. It took until the 1820s, when Augustin-Louis Cauchy reconstructed calculus on the foundation of limits, for Leibniz's geometric algebra to gain its modern rigor. The corner $\Delta u\,\Delta v$ is not discarded by arbitrary fiat; in the limit as $\Delta x \To 0$, its contribution $(\Delta u/ \Delta x)\Delta v$ vanishes because $\Delta v \To 0$. Leibniz's geometric intuition was completely vindicated; Cauchy supplied the rigorous backbone.
\end{historybox}

\section{The Chain Rule}
\label{ch11:chain-rule}

The chain rule is the most important differentiation rule of all. The power, product, and
quotient rules handle functions \emph{built by} arithmetic; the chain rule handles
functions built by \emph{composition} -- one function fed into another, like
$(3x+1)^{10}$, $\sqrt{x^2+1}$, $\sin(x^2)$, or $e^{-t/2}$. Almost every interesting
function in science is a composition, so almost every derivative you ever compute will
lean on this rule.

\subsection{The intuition: rates multiply}

Forget formulas for a moment and think about \emph{linked rates of change}. Suppose a
car's fuel use depends on its speed, and its speed depends on time. Ask: how fast is fuel
consumption changing with time?

Suppose that for each $1$ km/h of extra speed the car burns $0.2$
extra litres per hour (that is $\dd(\text{fuel})/\dd(\text{speed}) = 0.2$), and that right
now the car is accelerating at $3$ km/h per second
($\dd(\text{speed})/\dd t = 3$). Then per second the speed rises by $3$ km/h, and each of
those raises fuel use by $0.2$ litre/h, so fuel use climbs by $0.2 \times 3 = 0.6$ litre/h
per second. \emph{The rates multiply.}

\begin{intuitionbox}
\textbf{Gears.} Picture two meshed gears. If the first turns $3$ times for each turn of
the input shaft, and the second turns $2$ times for each turn of the first, then the
second turns $3\times 2 = 6$ times per turn of the input. Compose the linkages, multiply
the gear ratios. A derivative is precisely a ``how-much-output-per-unit-input ratio,''
and composing functions chains the linkages -- so you multiply the derivatives. That is
the entire content of the chain rule, before any symbols.
\end{intuitionbox}

In Leibniz notation the chain rule is almost self-evident. If $y$ depends on $u$ and $u$
depends on $x$, then
\[
  \dv{y}{x} = \dv{y}{u}\cdot\dv{u}{x},
\]
and it \emph{looks} as if the $\dd u$'s simply cancel like fractions. That visual
cancellation is the whole reason Leibniz's notation has survived three centuries -- it
makes the chain rule unforgettable. (It is not \emph{literally} fraction cancellation,
because $\dd y/\dd x$ is a limit, not a quotient of numbers; but it behaves enough like
one that the mnemonic never lies to you here.)

\subsection{Stating it in function notation}

Let $y = f(g(x))$ -- the ``outer'' function $f$ applied to the ``inner'' function $g$.
Setting $u = g(x)$, we have $y = f(u)$, and the chain rule reads:

\begin{theorem}[Chain rule]\label{ch11:thm-chain}
If $g$ is differentiable at $x$ and $f$ is differentiable at $g(x)$, then the composition
$f\circ g$ is differentiable at $x$, and
\[
  \ddx{f\big(g(x)\big)} = f'\big(g(x)\big)\cdot g'(x).
\]
In words: \emph{derivative of the outside (evaluated at the inside) times derivative of
the inside.}
\end{theorem}

\subsection{Why the naive proof almost works -- and how to fix it}

The natural proof writes the difference quotient and inserts a clever ``multiply by
$1$.'' Let $u=g(x)$ and let $\Delta u = g(x+h)-g(x)$ be the change in the inner function.
Then the change in $y$ is $\Delta y = f(g(x+h)) - f(g(x)) = f(u+\Delta u) - f(u)$, and we
write
\[
  \frac{\Delta y}{h}
  = \frac{f(u+\Delta u) - f(u)}{h}
  = \underbrace{\frac{f(u+\Delta u) - f(u)}{\Delta u}}_{\To\, f'(u)}
    \cdot
    \underbrace{\frac{\Delta u}{h}}_{\To\, g'(x)}.
\]
As $h\To 0$ we have $\Delta u \To 0$ (because $g$ is continuous), so the first factor
tends to $f'(u)$ and the second to $g'(x)$, giving $f'(u)\,g'(x)$. Beautiful -- and almost
correct.

\begin{trapbox}
The catch: we divided by $\Delta u$, but $\Delta u$ might be \emph{zero} for some small
$h$ (imagine $g$ momentarily flat). Then the fraction $\dfrac{f(u+\Delta u)-f(u)}{\Delta
u}$ is $0/0$, undefined, and the slick cancellation is illegal. This is not pedantic
fuss; it is a genuine gap that a rigorous proof must close.
\end{trapbox}

\subsection{A rigorous proof}

The standard repair, due in spirit to Carath\'eodory, replaces the troublesome difference
quotient by a single function that is well-behaved \emph{even when $\Delta u = 0$}. Define
the ``error-corrected slope'' of $f$ at $u$:
\[
  \phi(t) =
  \begin{cases}
    \dfrac{f(t)-f(u)}{t-u}, & t\ne u,\\[1.0em]
    f'(u), & t = u.
  \end{cases}
\]
Because $f$ is differentiable at $u$, the first branch tends to $f'(u)$ as $t\To u$, which
is exactly the value we assigned in the second branch -- so \emph{$\phi$ is continuous at
$u$}. Moreover, for \emph{every} $t$ (including $t=u$!), multiplying back across gives the
clean identity
\[
  f(t)-f(u) = \phi(t)\,(t-u).
\]
When $t\ne u$ this is clearing the fraction; when $t=u$ both sides are $0$.
By expressing the relationship as a product rather than a fraction, division by zero
is bypassed cleanly.

\begin{proof}[Proof of the chain rule]
Put $t = g(x+h)$ and recall $u = g(x)$, so $t - u = \Delta u = g(x+h)-g(x)$. The identity
$f(t)-f(u)=\phi(t)(t-u)$ becomes
\[
  f\big(g(x+h)\big) - f\big(g(x)\big) = \phi\big(g(x+h)\big)\cdot\big(g(x+h)-g(x)\big).
\]
Divide by $h$ (this is now perfectly safe -- no division by $\Delta u$ anywhere):
\[
  \frac{f(g(x+h)) - f(g(x))}{h}
  = \phi\big(g(x+h)\big)\cdot\frac{g(x+h)-g(x)}{h}.
\]
Now let $h\To 0$. Since $g$ is differentiable, it is continuous, so $g(x+h)\To g(x)=u$;
and since $\phi$ is continuous at $u$, $\phi(g(x+h))\To \phi(u) = f'(u)$. The second factor
tends to $g'(x)$ by definition. Therefore the product tends to
\[
  f'(u)\cdot g'(x) = f'\big(g(x)\big)\,g'(x),
\]
which is the chain rule.
\end{proof}

\begin{whybox}
Both proofs agree on the punchline -- ``the two rates
multiply'' -- and differ only in plumbing. The naive version divides by $\Delta u$ and
prays it is nonzero; the Carath\'eodory version replaces ``difference quotient of $f$''
with the genuinely continuous function $\phi$, which equals that quotient when $\Delta u
\ne 0$ but is also defined (as $f'(u)$) when $\Delta u = 0$. The product
$f(t)-f(u)=\phi(t)(t-u)$ holds with no exceptions, so we never divide by something that
might vanish. The moral is general: \emph{when a proof stumbles over a possible $0/0$,
look for a way to write the same quantity as a product instead of a quotient.}
\end{whybox}

\begin{figure}[H]
\centering
\begin{tikzpicture}[scale=0.92]
  % Top Lane: x-axis
  \node[anchor=west, font=\footnotesize\scshape, text=sourceblue!80!black] at (-0.5, 3.8) {Input Domain ($x$)};
  \draw[->, >=stealth, thick, gray!60] (-0.5, 3.2) -- (8.5, 3.2);
  \draw[very thick, sourceblue] (2.0, 3.2) -- (3.2, 3.2);
  \fill[sourceblue] (2.0, 3.2) circle (2pt) node[above, font=\small] {$x$};
  \fill[sourceblue] (3.2, 3.2) circle (2pt) node[above, font=\small] {$x + \Delta x$};
  \draw[<->, >=stealth, sourceblue!80!black] (2.0, 2.7) -- (3.2, 2.7) node[midway, below, font=\footnotesize] {step $\Delta x$};

  % Corridor 1 -> 2
  \draw[->, >=stealth, thick, sourceblue!70!black] (4.0, 2.4) -- (4.0, 1.6)
    node[midway, right, font=\small] {Mapping $g$: stretches by $g'(x) = \dfrac{\Delta u}{\Delta x}$};

  % Middle Lane: u-axis
  \node[anchor=west, font=\footnotesize\scshape, text=quotientorange!80!black] at (-0.5, 1.4) {Intermediate Value ($u = g(x)$)};
  \draw[->, >=stealth, thick, gray!60] (-0.5, 0.8) -- (8.5, 0.8);
  \draw[very thick, quotientorange] (1.5, 0.8) -- (4.2, 0.8);
  \fill[quotientorange] (1.5, 0.8) circle (2pt) node[above, font=\small] {$u$};
  \fill[quotientorange] (4.2, 0.8) circle (2pt) node[above, font=\small] {$u + \Delta u$};
  \draw[<->, >=stealth, quotientorange!80!black] (1.5, 0.3) -- (4.2, 0.3) node[midway, below, font=\footnotesize] {$\Delta u \approx g'(x)\,\Delta x$};

  % Corridor 2 -> 3
  \draw[->, >=stealth, thick, targetgreen!70!black] (4.0, 0.0) -- (4.0, -0.8)
    node[midway, right, font=\small] {Mapping $f$: stretches by $f'(u) = \dfrac{\Delta y}{\Delta u}$};

  % Bottom Lane: y-axis
  \node[anchor=west, font=\footnotesize\scshape, text=targetgreen!80!black] at (-0.5, -1.0) {Final Output ($y = f(u)$)};
  \draw[->, >=stealth, thick, gray!60] (-0.5, -1.6) -- (8.5, -1.6);
  \draw[very thick, targetgreen!70!black] (1.0, -1.6) -- (6.0, -1.6);
  \fill[targetgreen!70!black] (1.0, -1.6) circle (2pt) node[above, font=\small] {$y$};
  \fill[targetgreen!70!black] (6.0, -1.6) circle (2pt) node[above, font=\small] {$y + \Delta y$};
  \draw[<->, >=stealth, targetgreen!70!black] (1.0, -2.1) -- (6.0, -2.1) node[midway, below, font=\footnotesize] {$\Delta y \approx f'(u)\,\Delta u = f'(g(x))\,g'(x)\,\Delta x$};

  % Total rate callout box
  \draw[structurepurple!70, fill=structurepurple!10, rounded corners=2pt, thick] (9.0, 0.2) rectangle (13.6, 2.2);
  \node[font=\footnotesize\scshape, text=structurepurple!80!black] at (11.3, 1.8) {End-to-End Rate};
  \node[font=\small, align=center] at (11.3, 1.0) {$\dv{y}{x} = \dv{y}{u} \cdot \dv{u}{x}$\\[0.4em]$= f'\big(g(x)\big)\,g'(x)$};
\end{tikzpicture}
\caption{The chain rule as consecutive stretch factors across number lines. A small nudge
$\Delta x$ on the input axis is scaled by $g'(x)$ to produce the nudge $\Delta u$ on the intermediate
axis. That intermediate nudge is in turn scaled by $f'(u)$ on the output axis to produce $\Delta y$.
Because the transformations act in sequence, their local scaling factors multiply.}
\label{ch11:fig-chain}
\end{figure}

\subsection{Using it: peel the onion from the outside in}

In practice the chain rule is a peeling procedure. Identify the outermost operation,
differentiate it while leaving the inside untouched, then multiply by the derivative of
the inside -- and recurse if the inside is itself a composition.

\begin{example}[A power of a polynomial]\label{ch11:ex-chain1}
Differentiate $y = (3x+1)^{10}$.

Outer function: ``something to the $10$th,'' with derivative $10(\,\cdot\,)^9$. Inner
function: $3x+1$, with derivative $3$. So
\[
  y' = 10(3x+1)^{9}\cdot 3 = 30(3x+1)^{9}.
\]
Expanding $(3x+1)^{10}$ to differentiate term by term would mean eleven terms and a
binomial expansion. The chain rule does it in one line.
\end{example}

\begin{example}[A square root]\label{ch11:ex-chain2}
Differentiate $y = \sqrt{x^2+1} = (x^2+1)^{1/2}$.

Outer: $(\,\cdot\,)^{1/2}$, derivative $\tfrac12(\,\cdot\,)^{-1/2}$. Inner: $x^2+1$,
derivative $2x$. Thus
\[
  y' = \tfrac12 (x^2+1)^{-1/2}\cdot 2x = \frac{x}{\sqrt{x^2+1}} .
\]
Figure~\ref{ch11:fig-chain-graph} shows this slope as an actual tangent line.
This finally justifies the square-root step we borrowed in
Example~\ref{ch11:ex-prod2}.
\end{example}

\begin{figure}[H]
\centering
\begin{tikzpicture}
\begin{axis}[
  width=11cm, height=6.6cm,
  xlabel={$x$}, ylabel={$y$},
  domain=-3:3, samples=120,
  axis lines=middle,
  xmin=-3.2, xmax=3.2, ymin=0, ymax=3.6,
  legend style={at={(0.5,1.05)}, anchor=south, legend columns=2, font=\footnotesize},
  legend cell align=left,
  enlargelimits=false,
]
\addplot[curveblue, very thick] {sqrt(x^2+1)};
\addlegendentry{$y=\sqrt{x^2+1}$}
% tangent at x=2: y' = 2/sqrt(5), point (2, sqrt5)
\addplot[tangentred, thick, samples=2, domain=0.4:3] {(2/sqrt(5))*(x-2)+sqrt(5)};
\addlegendentry{tangent at $x=2$, slope $\tfrac{2}{\sqrt5}$}
\addplot[only marks, mark=*, black] coordinates {(2,2.2360679)};
\end{axis}
\end{tikzpicture}
\caption{The chain rule predicts the slope of a composition. For
$y=\sqrt{x^2+1}$ the chain rule gives $y'=x/\sqrt{x^2+1}$; at $x=2$ this is
$2/\sqrt5\approx 0.894$, the slope of the tangent line drawn here. The inner
derivative $2x$ is exactly what tilts the tangent away from the slope of the
bare square root.}
\label{ch11:fig-chain-graph}
\end{figure}

\begin{example}[The general power-chain shortcut]\label{ch11:ex-chain-genpow}
Combining the power rule with the chain rule gives a pattern you will use constantly:
\[
  \ddx{\big(g(x)\big)^n} = n\big(g(x)\big)^{n-1}\,g'(x).
\]
For instance $\ddx{(x^3 - 2x)^4} = 4(x^3-2x)^3(3x^2-2)$.
\end{example}

\begin{example}[Two layers deep]\label{ch11:ex-chain3}
Differentiate $y = \big(\sqrt{x}+1\big)^{3}$ -- a composition inside a composition.

Peel from the outside. Outermost: $(\,\cdot\,)^3$, derivative $3(\,\cdot\,)^2$. Its inside
is $\sqrt{x}+1$, whose derivative is $\tfrac{1}{2\sqrt{x}}$ (the $+1$ contributes nothing).
So
\[
  y' = 3\big(\sqrt{x}+1\big)^{2}\cdot\frac{1}{2\sqrt{x}}
     = \frac{3\big(\sqrt{x}+1\big)^{2}}{2\sqrt{x}} .
\]
If there had been a third layer, we would simply have multiplied by one more inner
derivative. The chain rule composes as many times as you need: it is gear ratios all the
way down.
\end{example}

\begin{trapbox}
The single most common chain-rule error is \emph{forgetting the inner derivative}. People
write $\ddx{(3x+1)^{10}} = 10(3x+1)^9$ and stop -- dropping the $\times 3$. The tell-tale
check: if the inside is just $x$, the inner derivative is $1$ and you are back to the
plain power rule; the moment the inside is anything more interesting, that $\times g'(x)$
factor is doing real work. \emph{Never let go of the inside.}
\end{trapbox}

\subsection{Combining all the rules}

Real problems braid the rules together. The skill is recognizing the \emph{outermost}
structure first -- is the whole expression a sum, a product, a quotient, or a composition?
-- and applying the matching rule, then recursing into the pieces.

\begin{example}[Product and chain together]\label{ch11:ex-combo1}
Differentiate $y = x^2\,(2x+3)^5$.

The outermost structure is a \emph{product}: $f = x^2$ times $g = (2x+3)^5$. Product rule
first:
\[
  y' = f'g + fg' = 2x\,(2x+3)^5 + x^2\cdot\ddx{(2x+3)^5}.
\]
The remaining derivative $\ddx{(2x+3)^5}$ needs the \emph{chain} rule:
$5(2x+3)^4\cdot 2 = 10(2x+3)^4$. Substitute back:
\[
  y' = 2x(2x+3)^5 + 10x^2(2x+3)^4 .
\]
You can factor out the common $2x(2x+3)^4$ to tidy up:
\[
  y' = 2x(2x+3)^4\big[(2x+3) + 5x\big] = 2x(2x+3)^4(7x+3).
\]
\end{example}

\begin{example}[Quotient and chain together]\label{ch11:ex-combo2}
Differentiate $y = \dfrac{(x+1)^3}{x^2+1}$.

Outermost structure: a \emph{quotient}. Top $f = (x+1)^3$ with $f' = 3(x+1)^2$ (chain
rule, inner derivative $1$); bottom $g = x^2+1$ with $g' = 2x$. Quotient rule:
\[
  y' = \frac{3(x+1)^2(x^2+1) - (x+1)^3(2x)}{(x^2+1)^2}.
\]
Factor $(x+1)^2$ from the numerator:
\[
  y' = \frac{(x+1)^2\big[3(x^2+1) - 2x(x+1)\big]}{(x^2+1)^2}
     = \frac{(x+1)^2\,(x^2 - 2x + 3)}{(x^2+1)^2}.
\]
(The bracket simplifies as $3x^2+3-2x^2-2x = x^2-2x+3$.)
\end{example}

\begin{walkbox}
On your next walk, try differentiating $\big((x^2+1)^3\big)$ \emph{two ways} in your head:
once by the chain rule directly, and once by first imagining it as the product
$(x^2+1)(x^2+1)(x^2+1)$ and using the three-factor product rule. Convince yourself both
give $6x(x^2+1)^2$. Then ask: why \emph{must} they agree? (Because they are the same
function -- and a function has only one derivative. The rules are different roads to one
destination, and the destination cannot move.)
\end{walkbox}

\begin{applicationbox}
\textbf{Where the chain rule earns its living.} In machine learning, training a neural
network means minimizing a loss that is a deep composition -- layer after layer after
layer of functions stacked on top of one another. \emph{Backpropagation}, the algorithm
that makes modern deep learning possible, is nothing but the chain rule applied
mechanically through that stack, multiplying the local rate of change at each layer to get
the overall rate of change of the loss with respect to each weight. When you read that a
network ``propagates gradients backward,'' it is computing $f'(g(x))\cdot g'(x)$ over and
over, exactly the rule you just proved. Newton's gear ratios are training the chatbots.
\end{applicationbox}

\begin{connectionbox}
The chain rule is the seed of multivariable calculus. When $y$ depends on several
intermediate quantities that each depend on $x$, the single product $f'(g(x))g'(x)$ blooms
into a \emph{sum} of such products -- the multivariable chain rule -- and that sum, written
compactly, is a dot product of the gradient with a velocity vector. The gradient,
directional derivatives, and the backpropagation above are all the chain rule wearing
different clothes. We will return to this when we reach functions of several variables.
\end{connectionbox}

\section{A Summary You Can Trust}
\label{ch11:summary}

You now hold the complete elementary toolkit. Every entry below was \emph{derived}, not
decreed, and each carries a one-line reason you can reconstruct.

\begin{center}
\renewcommand{\arraystretch}{1.5}
\begin{tabular}{@{}lll@{}}
\toprule
\textbf{Rule} & \textbf{Formula} & \textbf{Why} \\
\midrule
Constant & $\ddx{c} = 0$ & flat line, no change \\
Identity & $\ddx{x} = 1$ & change passes through \\
Power & $\ddx{x^n} = n x^{n-1}$ & linear term of $(x+h)^n$ survives \\
Constant multiple & $(kf)' = k f'$ & vertical stretch scales all slopes \\
Sum & $(f\pm g)' = f' \pm g'$ & rates add \\
Product & $(fg)' = f'g + fg'$ & rectangle grows on two sides \\
Quotient & $\left(\dfrac{f}{g}\right)' = \dfrac{f'g - fg'}{g^2}$ & top up, bottom down \\
Chain & $\big(f(g(x))\big)' = f'(g(x))\,g'(x)$ & rates multiply \\
\bottomrule
\end{tabular}
\end{center}

\begin{aestheticbox}
Step back and admire the architecture. Two atoms ($\ddx{c}=0$ and $\ddx{x}=1$), one
structural principle (linearity), and three combination rules (product, quotient, chain)
generate the derivative of \emph{every} function you can build from powers by the four
arithmetic operations and composition. That is an enormous reach from a tiny seed. And
none of it was arbitrary: each rule is forced by what ``rate of change'' \emph{means}
once you take the limit seriously. This is the recurring joy of calculus -- a handful of
plain ideas, pushed carefully, organize an entire universe of computation. In the next
chapter we add the derivatives of $e^x$, $\ln x$, and the trig functions, and with those
few new atoms the same machine differentiates essentially anything in the analyst's
toolbox.
\end{aestheticbox}

\subsection{The Scope and Horizon of Differentiation Rules}
\label{ch11:what-unlocked}

\begin{connectionbox}
Prior to Leibniz's 1684 operational calculus, finding the slope of a curve was an artisanal craft. Mathematicians had to devise custom geometric constructions or labor through bespoke difference-quotient expansions for every new curve. Leibniz's rules automated differentiation. By recognizing that functions are assembled algebraically from atomic pieces, calculus gained a grammar: any function constructed from powers, sums, products, quotients, and compositions could be differentiated mechanically on paper without returning to the limit definition from scratch.

Yet for all their algebraic power, differentiation rules provide purely \emph{local} information. The derivative $f'(x_0)$ reports the instantaneous velocity or the tilt of the tangent line at one isolated point $x_0$. It does not, on its own, reveal the global landscape: it cannot locate the absolute peak or deepest valley on a closed domain, it cannot guarantee that a function is monotonic or invertible across an interval merely from an algebraic formula, and it gives no automatic bound on how quickly a linear approximation deteriorates as you step away from $x_0$.

Turning these local rates into global understanding requires three subsequent bridges across the curriculum:
\begin{enumerate}
  \item In Chapter~14 (\emph{Linear Approximation and Differentials}), we use the local slope to construct linear models and quantify how differentials propagate measurement errors.
  \item In Chapter~15 (\emph{The Mean Value Theorem}), we prove the foundational theoretical bridge that allows local derivative signs ($f'(x) > 0$) to guarantee global behavior (the function strictly rises across an entire interval).
  \item In Chapter~16 (\emph{Maxima, Minima, and the Shape of a Graph}), we classify critical points and inflection points, translating local derivatives into complete qualitative portraits of a function's global geometry.
\end{enumerate}
\end{connectionbox}

\section{Exercises}
\label{ch11:exercises}

\begin{exercise}[Power-rule drill]
Differentiate each, rewriting roots and reciprocals as powers first.
\begin{enumerate}
  \item $f(x) = 6x^4 - 2x^3 + x - 11$.
  \item $g(x) = \dfrac{5}{x^3} - \dfrac{2}{\sqrt{x}} + 7$.
  \item $h(x) = \sqrt[3]{x^2} + 4x^{-5}$.
\end{enumerate}
\end{exercise}

\begin{exercise}[Product rule]
Differentiate using the product rule, then (where feasible) check by expanding first.
\begin{enumerate}
  \item $y = (2x^2 - 1)(x^3 + x)$.
  \item $y = x^3\sqrt{x+4}$.
  \item $y = (x^2+1)(x^2+2)(x^2+3)$ \;(use the three-factor product rule).
\end{enumerate}
\end{exercise}

\begin{exercise}[Quotient rule]
Differentiate and simplify.
\begin{enumerate}
  \item $y = \dfrac{x}{x^2+1}$.
  \item $y = \dfrac{2x-3}{4x+5}$.
  \item $y = \dfrac{x^2}{(x-1)^2}$ \;(quotient rule, with the chain rule on the bottom).
\end{enumerate}
\end{exercise}

\begin{exercise}[Chain rule]
Differentiate. Peel from the outside in, and do not drop the inner derivative.
\begin{enumerate}
  \item $y = (5x^2 - 3)^{7}$.
  \item $y = \sqrt{4 - x^2}$.
  \item $y = \big((x+1)^2 + 1\big)^{3}$ \;(two layers).
\end{enumerate}
\end{exercise}

\begin{exercise}[All rules braided together]
Differentiate. First identify the outermost structure (sum, product, quotient, or
composition).
\begin{enumerate}
  \item $y = x^2\sqrt{x^2+1}$.
  \item $y = \dfrac{(3x-1)^4}{x^2+1}$.
  \item $y = \big(x + \sqrt{x}\big)^5$.
\end{enumerate}
\end{exercise}

\begin{exercise}[Conceptual: spot the error]
A student writes $\ddx{\dfrac{x^2}{x+1}} = \dfrac{2x}{1} = 2x$ by ``differentiating top and
bottom separately.'' Explain in one or two sentences why this is wrong, name the rule that
should have been used, and give the correct derivative.
\end{exercise}

\begin{exercise}[Conceptual: why two terms?]
In your own words, explain why $\ddx{fg}$ has \emph{two} terms rather than one, using the
growing-rectangle picture. What does the third piece (the corner $\Delta f\,\Delta g$)
represent, and why does it not appear in the final formula?
\end{exercise}

\begin{exercise}[Reconstruct a rule]
Without looking back, derive the quotient rule from the product rule by writing
$f = q\,g$ where $q = f/g$, differentiating, and solving for $q'$. Then check your result
by computing $\ddx{1/x}$ two ways: from the quotient rule and from the limit definition.
\end{exercise}

\begin{exercise}[Proof: the power rule for negative integers]
Assume the power rule holds for positive integers and that the reciprocal rule holds. Let
$n$ be a positive integer and set $x^{-n} = 1/x^{n}$. Use the reciprocal rule
$\ddx{1/g} = -g'/g^2$ to prove
\[
  \ddx{x^{-n}} = -n\,x^{-n-1},
\]
thereby extending the power rule to negative integer exponents.
\end{exercise}

\begin{exercise}[Proof: a three-factor product rule]
Prove that $(fgh)' = f'gh + fg'h + fgh'$ by treating $fgh$ as the product of $f$ with the
single function $(gh)$ and applying the two-factor product rule twice. Then state, without
proof, the pattern for a product of $n$ factors.
\end{exercise}

\begin{exercise}[Proof: the chain rule cannot be skipped]
Let $F(x) = f(g(x))$ with $f(u) = u^2$ and $g(x) = x^2 + 1$.
\begin{enumerate}
  \item Compute $F'(x)$ by the chain rule.
  \item Compute $F'(x)$ instead by first expanding $F(x) = (x^2+1)^2 = x^4 + 2x^2 + 1$ and
        differentiating directly.
  \item Verify the answers agree, and explain what would have gone wrong had you mistakenly
        written $F'(x) = f'(g(x)) = 2(x^2+1)$ without the factor $g'(x)$.
\end{enumerate}
\end{exercise}

\begin{exercise}[A limit that hides a derivative]
Recognize each limit as a derivative $f'(a)$ for a suitable $f$ and $a$, and evaluate it
using the rules of this chapter rather than the limit directly.
\begin{enumerate}
  \item $\displaystyle\lim_{h\To 0}\frac{(2+h)^5 - 2^5}{h}$.
  \item $\displaystyle\lim_{h\To 0}\frac{\sqrt{9+h} - 3}{h}$.
\end{enumerate}
\end{exercise}
